% © 2026 Ing. David Jaroš — CC BY-NC-ND 4.0
% Licensed under CC BY-NC-ND 4.0. See LICENSE.md.
\documentclass[12pt,a4paper]{article}
\usepackage[a4paper,margin=2.5cm]{geometry}
\usepackage{amsmath,amssymb,amsthm,mathtools,mathrsfs}
\usepackage[most]{tcolorbox}
\usepackage[colorlinks=true,linkcolor=blue]{hyperref}
\newtheorem{definition}{Definition}[section]
\newtheorem{theorem}[definition]{Theorem}
\newtheorem{corollary}[definition]{Corollary}
\newtheorem{proposition}[definition]{Proposition}
\theoremstyle{remark}\newtheorem{remark}[definition]{Remark}
\newcommand{\B}{\mathbb B}
\title{GAP-10I: Flat Representers and the Connection-Representation Obstruction}
\author{Ing. David Jaro\v{s}}
\date{16 July 2026}
\begin{document}
\maketitle

\begin{abstract}
The UBT tetrad is constrained by
$E_\mu=\mathcal N_0^{-1/2}D_\mu\Theta$, so the local rank theorem alone does
not establish that a single $\Theta$ generates the required tetrads.  This
note closes two sharply defined subproblems.  First, every constant Lorentz
tetrad, and in particular Minkowski spacetime, has an explicit affine
$\Theta$ representer in the flat inertial sector; it also solves the standard
source-free second-order master operator because all second derivatives
vanish.  Second, a generic invertible $\Theta$ with a naive one-sided
connection cannot support nonzero curvature under torsion-free tetrad
compatibility.  The natural biquaternionic two-sided/bimodule derivative avoids
this flatness obstruction: its left and right curvatures need only be conjugate
through $\Theta$.  This does not solve curved-space integrability, but it
rigorously narrows the admissible canonical connection structure.
\end{abstract}

\section{Setup}
Let $\B=\mathbb C\otimes\mathbb H$ and let the Lorentz slice be
\begin{equation}
 W_L=\{i x^0\mathbf1+x^k\mathbf e_k\mid x^a\in\mathbb R\}.
\end{equation}
For a UBT field $\Theta$ define
\begin{equation}
 E_\mu:=\mathcal N_0^{-1/2}D_\mu\Theta,
 \qquad
 \frac12(E_\mu^\sharp E_\nu+E_\nu^\sharp E_\mu)
 =g_{\mu\nu}\mathbf1.
 \label{eq:tetrad-metric}
\end{equation}
The issue addressed here is existence and compatibility of $\Theta$, not the
already proved rank of $e_\mu{}^a\mapsto g_{\mu\nu}$.

\section{Exact flat representer}
\begin{theorem}[Affine $\Theta$ representer for a constant tetrad]
\label{thm:flat-representer}
Let $\bar E_\mu\in W_L$ be four constant, linearly independent
biquaternions.  In a flat inertial gauge with $\Omega_\mu=0$, define
\begin{equation}
 \boxed{
 \Theta_{\mathrm{aff}}(x)
 =\Theta_0+\sqrt{\mathcal N_0}\,\bar E_\mu x^\mu,}
 \label{eq:affine-theta}
\end{equation}
where $\Theta_0\in\B$ is constant.  Then
\begin{equation}
 \mathcal N_0^{-1/2}\partial_\mu\Theta_{\mathrm{aff}}
 =\bar E_\mu,
\end{equation}
and Eq.~\eqref{eq:tetrad-metric} yields the constant Lorentz metric
\begin{equation}
 \bar g_{\mu\nu}\mathbf1
 =\frac12(\bar E_\mu^\sharp\bar E_\nu
         +\bar E_\nu^\sharp\bar E_\mu).
\end{equation}
\end{theorem}

\begin{proof}
Differentiating Eq.~\eqref{eq:affine-theta} gives
$\partial_\mu\Theta_{\mathrm{aff}}
=\sqrt{\mathcal N_0}\bar E_\mu$.  The metric statement is the central
anticommutator identity on $W_L$.
\end{proof}

\begin{corollary}[Minkowski UBT vacuum representer]
Choose
\begin{equation}
 \bar E_0=i\mathbf1,
 \qquad \bar E_k=\mathbf e_k.
\end{equation}
Then
\begin{equation}
 \boxed{
 \Theta_{\mathrm{SR}}(x)
 =\Theta_0+\sqrt{\mathcal N_0}
  \left(i x^0\mathbf1+x^k\mathbf e_k\right)}
 \label{eq:theta-sr}
\end{equation}
induces $g_{\mu\nu}=\operatorname{diag}(-1,1,1,1)$.  All second spacetime
derivatives vanish.  Therefore any flat source-free master operator whose
principal part is the constant-coefficient covariant Laplacian or d'Alembertian
annihilates $\Theta_{\mathrm{SR}}$.
\end{corollary}

\begin{remark}
This is stronger than merely saying that one may set $\Omega_\mu=0$ in special
relativity: it gives an explicit single-field UBT configuration whose first
derivatives are the Minkowski tetrad.  It does not by itself fix boundary
conditions in compact imaginary time or prove uniqueness of the vacuum.
\end{remark}

\section{Why a naive one-sided connection is too restrictive}
Consider a one-sided connection on $\Theta$,
\begin{equation}
 D^L_\mu\Theta=\partial_\mu\Theta+A_\mu\Theta,
 \qquad
 F^L_{\mu\nu}=\partial_\mu A_\nu-\partial_\nu A_\mu+[A_\mu,A_\nu].
 \label{eq:left-connection}
\end{equation}
Then
\begin{equation}
 [D^L_\mu,D^L_\nu]\Theta=F^L_{\mu\nu}\Theta.
 \label{eq:left-curvature}
\end{equation}

\begin{theorem}[One-sided invertible-field flatness obstruction]
\label{thm:one-sided-no-go}
Assume on an open set that
\begin{enumerate}
 \item $E_\nu=\mathcal N_0^{-1/2}D^L_\nu\Theta$;
 \item the generated tetrad obeys torsion-free antisymmetric compatibility,
       $\mathscr D_{[\mu}E_{\nu]}=0$;
 \item the same one-sided connection is used in the induced differentiation
       of $D_\nu\Theta$;
 \item $\Theta$ is invertible as an element of
       $\B\simeq\operatorname{Mat}(2,\mathbb C)$.
\end{enumerate}
Then $F^L_{\mu\nu}=0$ on that open set.
\end{theorem}

\begin{proof}
Because the affine connection is torsion free, antisymmetrising the covariant
Hessian of the spacetime-scalar/internal field $\Theta$ gives
\begin{equation}
 \sqrt{\mathcal N_0}\,\mathscr D_{[\mu}E_{\nu]}
 =[D^L_\mu,D^L_\nu]\Theta
 =F^L_{\mu\nu}\Theta.
\end{equation}
The left-hand side vanishes by assumption, hence
$F^L_{\mu\nu}\Theta=0$.  Right multiplication by $\Theta^{-1}$ gives
$F^L_{\mu\nu}=0$.
\end{proof}

\begin{corollary}
A generic curved classical UBT branch cannot use the naive one-sided regular
representation together with invertible $\Theta$ and torsion-free tetrad
compatibility.  Noninvertible stabilizer solutions may evade the theorem but
are nongeneric and cannot be assumed to represent arbitrary GR geometries.
\end{corollary}

This theorem closes a route rather than the whole gap: the one-sided derivative
is not an acceptable generic curved-GR mechanism under the stated assumptions.

\section{The natural two-sided biquaternionic derivative}
Biquaternions form a bimodule over themselves, so left and right multiplication
are equally intrinsic.  Consider
\begin{equation}
 \boxed{
 D_\mu\Theta
 =\partial_\mu\Theta+A_\mu\Theta-\Theta B_\mu,}
 \label{eq:two-sided}
\end{equation}
with
\begin{align}
 F^A_{\mu\nu}
 &=\partial_\mu A_\nu-\partial_\nu A_\mu+[A_\mu,A_\nu],\\
 F^B_{\mu\nu}
 &=\partial_\mu B_\nu-\partial_\nu B_\mu+[B_\mu,B_\nu].
\end{align}
The relative minus sign is conventional and makes the right connection
transform in the standard bimodule form.

\begin{theorem}[Two-sided curvature identity]
\label{thm:two-sided-curvature}
For Eq.~\eqref{eq:two-sided},
\begin{equation}
 \boxed{
 [D_\mu,D_\nu]\Theta
 =F^A_{\mu\nu}\Theta-\Theta F^B_{\mu\nu}.}
 \label{eq:two-sided-curvature}
\end{equation}
\end{theorem}

\begin{proof}
Expand $D_\mu D_\nu\Theta-D_\nu D_\mu\Theta$.  Mixed terms containing one
ordinary derivative and one connection cancel pairwise, as do the terms
$A_\mu\Theta B_\nu$ after antisymmetrisation.  The remaining left products form
$F^A_{\mu\nu}\Theta$ and the remaining right products form
$-\Theta F^B_{\mu\nu}$.
\end{proof}

\begin{corollary}[Nonzero curvature is compatible with invertible $\Theta$]
Under torsion-free antisymmetric tetrad compatibility, the integrability
condition is
\begin{equation}
 F^A_{\mu\nu}\Theta=\Theta F^B_{\mu\nu}.
 \label{eq:curvature-intertwiner}
\end{equation}
If $\Theta$ is invertible, this becomes
\begin{equation}
 \boxed{F^A_{\mu\nu}=\Theta F^B_{\mu\nu}\Theta^{-1}.}
 \label{eq:curvature-conjugacy}
\end{equation}
Unlike the one-sided case, neither curvature must vanish.  The field $\Theta$
acts as an intertwiner between the left and right curvature representations.
\end{corollary}

\begin{remark}[What has and has not been selected]
The theorem does not yet derive $A_\mu$, $B_\mu$, their relation under the UBT
involutions, or their action from the canonical action.  It proves a structural constraint and exhibits a minimal algebra-native
escape from the one-sided obstruction: a two-sided/bimodule connection can
retain invertible $\Theta$, torsion-free compatibility, and nonzero
curvature simultaneously.  It is not proved to be the unique logically
possible escape.  A paired Lorentz spin connection is a natural candidate,
but its precise involution and normalization must be fixed in the canonical
action.
\end{remark}

\section{Implicit and possibly transcendental self-consistency}
Once the classical connection is reconstructed from the tetrad and torsion,
Eq.~\eqref{eq:two-sided} becomes schematically
\begin{equation}
 E_\mu=\mathcal N_0^{-1/2}
 \left[
 \partial_\mu\Theta
 +A_\mu[E,T]\Theta
 -\Theta B_\mu[E,T]
 \right].
 \label{eq:fixed-point}
\end{equation}
This is an implicit nonlinear first-order PDE or fixed-point system because
$E$ appears on both sides through the connection.  It is not automatically a
transcendental equation in the strict mathematical sense.  However, when
$\Theta(q,\tau)$ is restricted to a Jacobi-theta or other transcendental
function class, the complete UBT system may be both implicit and
transcendental.  These are distinct properties and should be stated
separately in formal proofs.

\section{Revised partial closure of GAP-10I}
\begin{tcolorbox}[colback=green!6!white,colframe=green!55!black,title=Defensible status]
\textbf{GAP-10I-SR: CLOSED [L1].}  Every constant Lorentz tetrad has the
explicit affine representer Eq.~\eqref{eq:affine-theta}; Minkowski spacetime is
therefore generated by a single $\Theta$ in the flat inertial branch.

\medskip
\textbf{GAP-10I-1S: CLOSED AS NO-GO [L1].}  A naive one-sided regular
connection with invertible $\Theta$ cannot support generic nonzero curvature
under torsion-free tetrad compatibility.

\medskip
\textbf{GAP-10I-2S: NARROWED [L1].}  The two-sided bimodule derivative avoids
the flatness obstruction and reduces integrability to the curvature-intertwiner
condition Eq.~\eqref{eq:curvature-intertwiner}.  This is a necessary structural
advance, not a proof of existence for arbitrary curved tetrads.

\medskip
\textbf{GAP-10I-PRESCRIBED: CLOSED [L1].}  For specified smooth
$(E,A,B)$ the exact solution and path-independence problem is the augmented
holonomy criterion proved in
\texttt{gap\_10i\_augmented\_holonomy.tex}.

\medskip
\textbf{GAP-10I-CURVED: NARROWED.}  Self-consistent action-level selection,
regularity, and global continuation for Schwarzschild, Kerr, FRW, and generic
Einstein geometries remain to be proved from the canonical UBT action.
\end{tcolorbox}

\end{document}
